\documentclass[11pt,a4paper]{article}

\usepackage[UTF8]{ctex}
\usepackage{booktabs}
\usepackage{enumitem}
\usepackage{fancyhdr}
\usepackage{geometry}
\usepackage{hyperref}
\usepackage{listings}
\usepackage[most]{tcolorbox}
\usepackage{xcolor}

\geometry{left=2.15cm,right=2.15cm,top=2.05cm,bottom=2.05cm}
\setlength{\headheight}{14pt}
\setlength{\parindent}{2em}
\setlength{\parskip}{0.28em}
\setlist[itemize]{leftmargin=2em,itemsep=0.25em,topsep=0.3em}
\setlist[enumerate]{leftmargin=2.25em,itemsep=0.25em,topsep=0.3em}

\definecolor{courseblue}{HTML}{2457A6}
\definecolor{coursegreen}{HTML}{18794E}
\definecolor{codered}{HTML}{A13D2D}
\definecolor{codegray}{HTML}{F4F6F8}
\definecolor{windowsblue}{HTML}{0067B8}
\definecolor{macgray}{HTML}{4B5563}
\hypersetup{colorlinks=true,linkcolor=courseblue,urlcolor=courseblue}

\lstdefinestyle{coursecode}{
  basicstyle=\ttfamily\small,
  backgroundcolor=\color{codegray},
  frame=single,
  rulecolor=\color{black!12},
  breaklines=true,
  columns=fullflexible,
  keepspaces=true,
  showstringspaces=false,
  keywordstyle=\color{courseblue}\bfseries,
  commentstyle=\color{coursegreen},
  stringstyle=\color{codered},
  aboveskip=0.55em,
  belowskip=0.55em
}
\lstset{style=coursecode}

\newtcolorbox{importantbox}{
  colback=courseblue!4,
  colframe=courseblue!65,
  boxrule=0.7pt,
  arc=1.5mm,
  left=1.5mm,right=1.5mm,top=1.2mm,bottom=1.2mm
}
\newtcolorbox{tipbox}{
  colback=coursegreen!4,
  colframe=coursegreen!65,
  boxrule=0.7pt,
  arc=1.5mm,
  left=1.5mm,right=1.5mm,top=1.2mm,bottom=1.2mm
}
\newcommand{\windows}{\subsection*{\textcolor{windowsblue}{Windows}}}
\newcommand{\macos}{\subsection*{\textcolor{macgray}{macOS}}}

\pagestyle{fancy}
\fancyhf{}
\lhead{AI3002 人工智能与机器学习基础}
\rhead{LAB0：开发环境配置}
\cfoot{\thepage}
\renewcommand{\headrulewidth}{0.4pt}
\fancypagestyle{plain}{%
  \fancyhf{}%
  \lhead{AI3002 人工智能与机器学习基础}%
  \rhead{LAB0：开发环境配置}%
  \cfoot{\thepage}%
  \renewcommand{\headrulewidth}{0.4pt}%
}

\title{\vspace{-1.35cm}\textbf{人工智能与机器学习基础\\LAB0：机器学习开发环境配置}}
\author{中国科学技术大学}
\date{2026 年秋季学期\quad 发布日期与截止日期待定}

\begin{document}

\maketitle
\pagestyle{fancy}
\thispagestyle{fancy}

\section*{实验目标与提交要求}

本实验帮助大家建立后续实验统一使用的 Python 开发环境。完成后，你应当能够使用 Git 获取课程代码，使用 Conda 管理独立环境，在 Visual Studio Code（VS Code）中选择正确的 Python 解释器，并正常导入 NumPy、scikit-learn 与 PyTorch。

本文分别给出 \textbf{Windows 10/11（64 位）}和 \textbf{macOS 11 或更高版本（Apple Silicon 或 Intel）}的完整操作。请选择自己的平台，按对应步骤从头完成；不要把两个平台的命令混在同一个终端中执行。

\begin{importantbox}
预计用时为 \textbf{60--90 分钟}。本实验必须单人完成，不要求 GPU，不要求撰写实验报告。最终只需向 OJ 提交验证脚本生成的 \textbf{64 位 SHA256 字符串}。截止前可以重复提交。
\end{importantbox}

实验文件目录如下：

\begin{lstlisting}
LAB0/
|-- environment.yml
|-- lab0.pdf
`-- verify_env.py
\end{lstlisting}

请不要修改 \texttt{verify\_env.py}。如果程序没有生成 SHA256，应根据错误提示修复环境，而不是删除检查代码。

对作业或实验内容有疑问时，请先搜索并通过课程仓库的 GitHub Issues 提问：
\begin{center}
  \url{https://github.com/Intelligent114/ALML_public/issues}
\end{center}
不要在 Issue 中发布本人学号、邮箱验证码、密码、SHA256 提交值或未公开答案。

\section{安装与配置 Git}

Git 是一个分布式版本控制系统。本课程使用 Git 发布和更新实验代码。

\windows

打开 Git 官方 Windows 安装页面：
\begin{center}
  \url{https://git-scm.com/install/windows}
\end{center}
下载并运行适合机器架构的 Git for Windows 安装器，保留默认选项即可。也可以在 PowerShell 中使用 Windows Package Manager：

\begin{lstlisting}[language=bash]
winget install --id Git.Git -e --source winget
\end{lstlisting}

安装完成后关闭并重新打开 PowerShell，运行：

\begin{lstlisting}[language=bash]
git --version
git config --global user.name "Your Name"
git config --global user.email "your_email@example.com"
git config --list --show-origin
\end{lstlisting}

在 PowerShell 中克隆课程仓库：

\begin{lstlisting}[language=bash]
cd $HOME
git clone https://github.com/Intelligent114/ALML_public.git
cd ALML_public
git status
git log -1
\end{lstlisting}

\macos

打开 Terminal（终端）。macOS 可以通过 Xcode Command Line Tools 安装 Apple 提供的 Git：

\begin{lstlisting}[language=bash]
xcode-select --install
\end{lstlisting}

在系统弹窗中选择安装，完成后关闭并重新打开 Terminal，再运行：

\begin{lstlisting}[language=bash]
git --version
git config --global user.name "Your Name"
git config --global user.email "your_email@example.com"
git config --list --show-origin
\end{lstlisting}

如果已经使用 Homebrew，也可以运行 \texttt{brew install git}；本实验不要求专门安装 Homebrew。在 Terminal 中克隆课程仓库：

\begin{lstlisting}[language=bash]
cd "$HOME"
git clone https://github.com/Intelligent114/ALML_public.git
cd ALML_public
git status
git log -1
\end{lstlisting}

\begin{tipbox}
Windows 与 macOS 后续获取更新的命令都是 \texttt{git pull}。如果出现本地修改冲突，请先备份自己的文件，再在 GitHub Issues 中说明平台、文件名和完整错误，不要盲目覆盖。
\end{tipbox}

\section{安装 Miniforge 并创建课程环境}

Conda 可以为不同项目建立相互隔离的 Python 环境。本课程使用由 conda-forge 维护的 Miniforge：
\begin{center}
  \url{https://github.com/conda-forge/miniforge/releases/latest}
\end{center}

\windows

下载 \texttt{Miniforge3-Windows-x86\_64.exe} 并运行。建议保留“创建开始菜单快捷方式”，不要勾选把 Miniforge 添加到系统 PATH，以免与其他 Python 冲突。

安装后从开始菜单打开 \textbf{Miniforge Prompt}，执行：

\begin{lstlisting}[language=bash]
conda --version
conda init powershell
\end{lstlisting}

关闭所有 Miniforge Prompt 和 PowerShell 窗口，再重新打开 PowerShell。进入 LAB0 目录并创建课程环境：

\begin{lstlisting}[language=bash]
cd $HOME\ALML_public\LABs\LAB0
conda env create -f environment.yml
conda activate ai3002
python --version
where.exe python
python -c "import sys; print(sys.executable)"
\end{lstlisting}

\macos

在 Apple 菜单的“关于本机”中查看芯片类型，或者在 Terminal 运行：

\begin{lstlisting}[language=bash]
uname -m
\end{lstlisting}

根据输出选择安装器：
\begin{itemize}
  \item \texttt{arm64}：下载 Apple Silicon 的 \path{Miniforge3-MacOSX-arm64.pkg}；
  \item \texttt{x86\_64}：下载 Intel 的 \path{Miniforge3-MacOSX-x86_64.pkg}。
\end{itemize}
运行签名的 PKG 安装器，并在安装选项中允许初始化当前 Shell。

安装后关闭并重新打开 Terminal，运行：

\begin{lstlisting}[language=bash]
conda --version
\end{lstlisting}

如果提示找不到 \texttt{conda}，执行下面的命令，然后再次关闭并重新打开 Terminal：

\begin{lstlisting}[language=bash]
~/miniforge3/bin/conda init zsh
\end{lstlisting}

进入 LAB0 目录并创建课程环境：

\begin{lstlisting}[language=bash]
cd "$HOME/ALML_public/LABs/LAB0"
conda env create -f environment.yml
conda activate ai3002
python --version
which python
python -c "import sys; print(sys.executable)"
\end{lstlisting}

两个平台都应看到 Python 3.11，解释器路径中应包含 \texttt{ai3002}。查看、退出或删除环境的命令在 Windows PowerShell 与 macOS Terminal 中相同：

\begin{lstlisting}[language=bash]
conda info --envs
conda deactivate
conda env remove -n ai3002
\end{lstlisting}

只有环境已经损坏且普通排错无法解决时才需要删除重建。后续每次做实验前，都应先运行 \texttt{conda activate ai3002}。

\section{安装 VS Code 并选择正确解释器}

VS Code 下载地址为 \url{https://code.visualstudio.com/}。安装后，在扩展面板安装 Microsoft 发布的 \textbf{Python} 扩展；Python Environments 扩展会协助发现和切换 Conda 环境。

\windows

运行 Windows 安装器时，建议启用“添加到 PATH”和“通过 Code 打开”选项。也可以在 PowerShell 中安装：

\begin{lstlisting}[language=bash]
winget install -e --id Microsoft.VisualStudioCode
\end{lstlisting}

重新打开 PowerShell，用 VS Code 打开整个课程仓库：

\begin{lstlisting}[language=bash]
code $HOME\ALML_public
\end{lstlisting}

按 \texttt{Ctrl+Shift+P}，选择：

\begin{center}
\texttt{Python: Select Interpreter} $\longrightarrow$ \texttt{ai3002 (Python 3.11)}
\end{center}

新建 VS Code 集成终端，确认提示符前有 \texttt{(ai3002)}，然后运行：

\begin{lstlisting}[language=bash]
where.exe python
python -c "import sys; print(sys.executable)"
\end{lstlisting}

\macos

下载 macOS 版本，将 Visual Studio Code 拖入 Applications。打开 VS Code，按 \texttt{Cmd+Shift+P}；如需在 Terminal 使用 \texttt{code}，先执行 \texttt{Shell Command: Install 'code' command in PATH}。

在 Terminal 中打开整个课程仓库：

\begin{lstlisting}[language=bash]
code "$HOME/ALML_public"
\end{lstlisting}

按 \texttt{Cmd+Shift+P}，选择：

\begin{center}
\texttt{Python: Select Interpreter} $\longrightarrow$ \texttt{ai3002 (Python 3.11)}
\end{center}

新建 VS Code 集成终端，确认提示符前有 \texttt{(ai3002)}，然后运行：

\begin{lstlisting}[language=bash]
which python
python -c "import sys; print(sys.executable)"
\end{lstlisting}

\begin{importantbox}
两个平台最常见的问题都不是“库没有安装”，而是 VS Code 选择了另一个 Python。Windows 用 \texttt{where.exe python}，macOS 用 \texttt{which python}；输出路径必须对应 \texttt{ai3002} 环境。
\end{importantbox}

\section{检查机器学习库}

\texttt{environment.yml} 会安装 NumPy、scikit-learn 与 PyTorch。LAB0 只检查 CPU 整数计算，不要求 CUDA，也不要求 Apple Metal/MPS。

\windows

在 VS Code 的 PowerShell 集成终端中运行：

\begin{lstlisting}[language=bash]
conda activate ai3002
cd $HOME\ALML_public\LABs\LAB0
python -c "import numpy, sklearn, torch; print('imports ok')"
python -c "import torch; print(torch.__version__)"
\end{lstlisting}

\macos

在 VS Code 的 zsh 集成终端中运行：

\begin{lstlisting}[language=bash]
conda activate ai3002
cd "$HOME/ALML_public/LABs/LAB0"
python -c "import numpy, sklearn, torch; print('imports ok')"
python -c "import torch; print(torch.__version__)"
\end{lstlisting}

Windows 没有 NVIDIA GPU 时，\texttt{torch.cuda.is\_available()} 返回 \texttt{False} 是正常的；macOS 通常也会返回 \texttt{False}。这不影响实验通过，不要为了 LAB0 额外安装 CUDA Toolkit 或修改 PyTorch 包。

\section{生成环境验证 SHA256}

验证脚本会检查 Python、Conda、Git 和三个机器学习库，并将确定性的整数结果与本人学号组合后计算 SHA256。

\windows

在 VS Code 的 PowerShell 集成终端中运行：

\begin{lstlisting}[language=bash]
cd $HOME\ALML_public\LABs\LAB0
conda activate ai3002
python verify_env.py --self-test
python verify_env.py --student-id PBXXXXXXXX
\end{lstlisting}

\macos

在 VS Code 的 zsh 集成终端中运行：

\begin{lstlisting}[language=bash]
cd "$HOME/ALML_public/LABs/LAB0"
conda activate ai3002
python verify_env.py --self-test
python verify_env.py --student-id PBXXXXXXXX
\end{lstlisting}

两个平台正常输出的末尾都类似：

\begin{lstlisting}
[PASS] Student ID: PBXXXXXXXX
[PASS] Protocol: AI3002-2026F-LAB0-v1

请将下面一行的 64 位字符串提交到 OJ：

0db3...中间省略...7a4c
\end{lstlisting}

最终 SHA256 只由协议版本、规范化后的学号和固定整数结果组成。软件版本、操作系统、安装路径、GPU 状态和 Git 姓名邮箱均不会上传，也不会影响结果。学号字母会统一转换为大写，最终 SHA256 为小写。

\section{提交到 OJ}

课程 OJ 地址为 \url{https://oj.temaurinum.moe/}。

\windows

使用 Edge、Chrome 或 Firefox 打开 OJ，登录后选择“LAB0：机器学习开发环境配置”。确认页面显示的学号正确，将 PowerShell 中生成的 64 位 SHA256 粘贴到输入框并提交。

\macos

使用 Safari、Chrome 或 Firefox 打开 OJ，登录后选择“LAB0：机器学习开发环境配置”。确认页面显示的学号正确，将 Terminal 中生成的 64 位 SHA256 粘贴到输入框并提交。

OJ 会使用当前登录账户中的学号重新计算期望结果，因此两个平台都应注意：

\begin{itemize}
  \item 不需要、也不能在 OJ 中另行填写学号；
  \item 复制其他同学的 SHA256 无法通过；
  \item 如果注册账户时填错学号，请在 GitHub Issue 中只说明“账户学号需要更正”，不要公开具体学号；
  \item 截止前可以重复提交，以出现“环境验证通过”为完成标志。
\end{itemize}

\begin{importantbox}
SHA256 是本实验的完成凭证，用于发现常见环境配置错误；它不是可信硬件证明。请独立完成安装过程，因为后续正式实验会直接依赖这个环境。
\end{importantbox}

\section{通过 GitHub Issues 提问}

所有作业与实验问题统一发布到：
\begin{center}
  \url{https://github.com/Intelligent114/ALML_public/issues/new/choose}
\end{center}

\windows

可以在浏览器中打开上面的地址，或在 PowerShell 中运行：

\begin{lstlisting}[language=bash]
start https://github.com/Intelligent114/ALML_public/issues/new/choose
\end{lstlisting}

\macos

可以在浏览器中打开上面的地址，或在 Terminal 中运行：

\begin{lstlisting}[language=bash]
open https://github.com/Intelligent114/ALML_public/issues/new/choose
\end{lstlisting}

提问前先搜索是否已有相同问题。新 Issue 应写明 HW/LAB 编号、Windows 或 macOS、系统版本、CPU 架构、执行的命令、完整错误文本和已经尝试的方法。请使用文字或代码块粘贴错误，避免只发一张无法搜索的截图。

\begin{importantbox}
GitHub Issue 是公开页面。不得发布真实学号、个人邮箱、验证码、密码、Cookie、SHA256 提交值、未公开答案或完整个人作业。需要私下处理的账户信息由助教在 Issue 中另行说明安全渠道。
\end{importantbox}

\section{常见问题与双平台排错}

\subsection*{找不到 git}

\windows
关闭并重新打开 PowerShell，运行 \texttt{where.exe git}。如果仍找不到，重新运行 Git for Windows 安装器，或执行 \texttt{winget install --id Git.Git -e --source winget}。

\macos
在 Terminal 运行 \texttt{xcode-select --install}，完成系统弹窗中的安装后重新打开 Terminal，再运行 \texttt{which git}。

\subsection*{找不到 conda}

\windows
从开始菜单打开 Miniforge Prompt，运行 \texttt{conda init powershell}，然后关闭所有终端并重新打开 PowerShell。

\macos
在 Terminal 运行 \texttt{~/miniforge3/bin/conda init zsh}，然后关闭并重新打开 Terminal。若 Miniforge 安装在其他目录，请在 GitHub Issue 中附上安装路径和 \texttt{ls \$HOME} 的相关输出，不要公开用户名目录中的私人文件。

\subsection*{VS Code 没有使用 ai3002}

\windows
按 \texttt{Ctrl+Shift+P} 重新运行 \texttt{Python: Select Interpreter}，选择 ai3002；关闭旧集成终端并新建终端，用 \texttt{where.exe python} 核对。

\macos
按 \texttt{Cmd+Shift+P} 重新运行 \texttt{Python: Select Interpreter}，选择 ai3002；关闭旧集成终端并新建终端，用 \texttt{which python} 核对。

\subsection*{可以导入 NumPy，但不能导入 sklearn 或 torch}

\windows

\begin{lstlisting}[language=bash]
cd $HOME\ALML_public\LABs\LAB0
conda activate ai3002
conda env update -n ai3002 -f environment.yml --prune
\end{lstlisting}

\macos

\begin{lstlisting}[language=bash]
cd "$HOME/ALML_public/LABs/LAB0"
conda activate ai3002
conda env update -n ai3002 -f environment.yml --prune
\end{lstlisting}

\subsection*{OJ 提示 SHA256 不匹配}

\windows
在 PowerShell 中先运行 \texttt{git pull}，再运行 \texttt{python verify\_env.py --self-test} 和本人学号命令；确认所有检查均为 PASS。

\macos
在 Terminal 中先运行 \texttt{git pull}，再运行 \texttt{python verify\_env.py --self-test} 和本人学号命令；确认所有检查均为 PASS。

如果仍未解决，请按上一节要求创建 GitHub Issue，但不要粘贴最终 SHA256 或本人学号。

\section{完成前自检清单}

提交前逐项核对；若某一项失败，先回到对应章节排查，再创建 GitHub Issue。

\windows

\begin{enumerate}
  \item \texttt{git --version} 和 \texttt{where.exe git} 均能正常输出；
  \item \texttt{conda info --envs} 中 ai3002 前有星号，\texttt{where.exe python} 的第一项属于 ai3002；
  \item 在仓库的 \path{LABs/LAB0} 目录运行库版本检查时，NumPy、scikit-learn 和 PyTorch 均能导入；
  \item \texttt{python verify\_env.py --self-test} 的全部检查均为 PASS；
  \item 使用本人学号生成 SHA256，且 OJ 显示“环境验证通过”。
\end{enumerate}

\macos

\begin{enumerate}
  \item \texttt{git --version} 和 \texttt{which git} 均能正常输出；
  \item \texttt{conda info --envs} 中 ai3002 前有星号，\texttt{which python} 指向 ai3002；
  \item 在仓库的 \path{LABs/LAB0} 目录运行库版本检查时，NumPy、scikit-learn 和 PyTorch 均能导入；
  \item \texttt{python verify\_env.py --self-test} 的全部检查均为 PASS；
  \item 使用本人学号生成 SHA256，且 OJ 显示“环境验证通过”。
\end{enumerate}

完成以上步骤后，请保留本地 ai3002 环境和课程仓库。后续实验只需进入仓库、执行 \texttt{git pull}、激活 ai3002，并在 VS Code 中确认解释器仍为该环境。

\vfill
\begin{center}
  \textcolor{gray}{--- End of LAB0 ---}
\end{center}

\end{document}
